\chapter{Závěr}
\label{sec:zaver}

Tento projekt představuje výsledek mé dlouhodobé práce na \emph{Surreal Bowling: Salvation Path}, který vznikl v rámci tří spolupracujících maturitních prací.

Během tvorby této práce se mi podařilo vytvořit komplexní audiovizuální systém pro 2D videohru, který zahrnuje pokročilé vizuální efekty pomocí Unity Universal Render Pipeline, profesionální zvukový design prostřednictvím FMOD Studio a vlastní hudbu v FL Studio, a moderní uživatelské rozhraní založené na UI Toolkitu. Práce mi umožnila prohloubit mé zkušenosti s implementací audia do herního kódu a efektivně propojit audio systém s herními mechanikami.

Mezi dovednosti, které jsem během práce rozvinul, patří práce s globálními objekty v Unity, správa zvukového mixu ve FMOD Studiu, tvorba hudebních kompozic v FL Studio a efektivní komunikace v týmu. FL Studio jsem používal již před touto prací, ale tato práce mi umožnila využít tyto dovednosti v profesionálnějším kontextu a propojit je s herním vývojem.

Spolupráce v týmu Still Thinking Studio byla klíčovým prvkem celého projektu. Naše týmová spolupráce se mohla díky této práci ověřit v praxi -- naučili jsme se efektivně komunikovat nápady, spolupracovat na společných cílech a respektovat role jednotlivých členů. Zkušenost s koordinací tří maturitních prací na jednom projektu nám ukázala, jak důležitá je kvalitní komunikace pro úspěšný herní projekt.

Je důležité zmínit, že projekt \emph{Surreal Bowling: Salvation Path} vznikl na game jamu od Game Devs Ostrava, kde jsme z 12 týmů zvítězili na prvním místě. Od té doby jsme se účastnili dalších soutěží a konferencí. Tato maturitní práce tedy není prvním krokem naší cesty, ale spíše příležitostí, jak na našem projektu dále pracovat a rozvíjet ho.

V současnosti se náš tým připravuje na účast na konferenci \emph{Game Access Conference '26} v Brně, kde budeme prezentovat nejen \emph{Surreal Bowling: Salvation Path}, ale také náš druhý projekt \emph{Dissonantia} -- top-down view roguelike hru. Tato konference představuje významný milník v naší developerské dráze a příležitost získat zpětnou vazbu od profesionálů v oboru.

Do budoucna plánujeme pokračovat ve vývoji obou her s cílem je vydat na platformě Steam. Věříme, že kvalita naší práce může zaujmout hráče a potenciálně i sponzory nebo vydavatele, kteří by mohli podpořit další rozvoj našich projektů.

Celkově tato práce ověřila naše schopnosti vyvíjet hry a díky příležitosti od naší školy jsme mohli pokračovat v práci na vlastním projektu. Získané zkušenosti budou cenným základem pro naši budoucnost v herním průmyslu.